\newpage

Notation:
\begin{itemize}
	\item $U\in\Uc$
	\item suff. stat. $X=t(U)$
	\item $\nu$ the underlying measure on $\Uc$
\end{itemize}

Assume $U$ some regular $d$-dim exp. family, represented by $\Pc$.
We write the density as
\begin{equation*}
	P(u) = \frac{1}{Z(\beta;\mu_*)} \exp(\beta^T X) p_{0;\mu^*}(u) \nu(du)
\end{equation*}

The mean value space of $\Pc$ is
\begin{equation*}
	\Mt_\Pc = \{\mu : \ERWi{P}{X} = \mu\} \qquad \text{for some } P \in \Pc
\end{equation*}

The canonical space is
\begin{equation*}
	\Bt = \{\beta : Z(\beta) = \int \exp(\beta^T X) p_{0;\mu_*}(u) \nu(du) < \infty \}
\end{equation*}

Gradient of the log partioning function gives the mean of the suff. stat. $X$.


\section{Problems in simple case}



\subsection{Why there is a mean in the cumulant function $Z$?}

TODO

\subsection{What is the mean of the carrying density}

$p_{\mu^*} = p_{0;\mu^*}$ is probably the most confusing part for me for the longest time.
But it is basically another way to state that
\begin{quote}
	When canonical parameter $\beta$ is set to zero,
	we got carrying density back.
\end{quote}

Then $\mu_*$ is the expectation of the mean of the suff. stat. $X=t(U)$
in the carrying density. 

At $\beta=0$,
\begin{equation*}
	P_{0;\mu_*}(u) = \frac{p_{0;\mu_*}(u)}{Z(0)} \nu(\dd{u}).
\end{equation*}.

Then there is a mean $\mu_*$ corresponding to that!
Right now, I would just write it explicitly.

\subsection{Why change of measure?}

Because we can write $\Pc$ with a different base measure instead of $\nu$.
We use $p_{\mu_1}$, one then can write the same set of distribution
\begin{equation*}
	\Pc_{\nu} = \Pc_{p_{\mu_1}}
\end{equation*}

The mean-value space after reparameterization stayed fixed
the suff. stat. and the distribution stays fixed, so is the expectation.

However, the original density reads like
\begin{align*}
	p_{\beta,\nu} &= \frac{1}{Z(\beta)} \exp\left(\beta^T t(u)\right) \nu(\dd{u}) \\
	&= \frac{1}{Z(\beta;\mu^*)} \exp\left(\beta^T t(u)\right) \nu(\dd{u})
\end{align*}

Instead of $\nu$, using a different carrying density allows us to re-paremeterize the
same set of distribution on the \emph{mean-value} space $\Mt_\Pc$ 
by changing the canonical spaces $\Bt_{\Pc}$

\subsection{Why generate alternative distribution?}

Starting with some distribution $Q$ on $\Uc$, density is $q(u)$.
We could write a set of distributions using the same
\begin{equation*}
	q_{\beta;\mu_*}(u) = \frac{1}{Z(\beta;\mu_*)} \exp(\beta^T t(u)) q(u)
\end{equation*}

$q(u)$ is what we originally have, $q_{\beta;\mu_*}$ is the \emph{generated} set of 
distributions on $\Uc$ by taking $q(u)$ as the carrying density.

WHY? we want two sets of distributions sharing the same set of suff. stat.

Why the mean $\mu_*$ appears in the cumulant function. 

\subsection{Benefits of sharing suff. stat.}

Because density $q_{\beta;\mu_*}(u)$ now shares the same suff. stat. as the null $\Pc$
If we define a function like (2.1 in simple), we can calculate the log of the ratio
\begin{equation*}
	f(\beta;\mu^*) = \log \ERWi{P_{\beta;\mu_*}}{\frac{q}{p}}
\end{equation*}

\subsection{Curvature}


%Simple case:
%
%1. pµ∗ = p0,µ∗ this is by direct calculation
%
%2. Since Q is assumed to have a moment generating function, the canonical domain Bq,µ∗ is
%nonempty and contains a neighborhood of 0
%
%Do we assume sharing same sufficient statistics?
%
%We assume throughout this paper that, for any considered alternative Q, there exists a µ∗∈Mp
%such that EX∼Q[X] = µ∗
